\subsection{Objective and experimental approach}

The first phase establishes the transmission and reception chain required for
the later direction-of-arrival experiments. Its purpose is to generate
repeatable radio signals, acquire their complex in-phase and quadrature
samples with an SDR (I/Q), and inspect or recover information from those samples.

\subsection{Preliminary laboratory experiments}

In the first laboratory exercise, an ESP32-C3 is configured to broadcast
packets using ESP-NOW on Wi-Fi channel 1, whose centre frequency is
\(2412~\mathrm{MHz}\). The transmitter selects a PHY rate of \(6~\mathrm{Mbit/s}\)
and broadcasts a packet every \(250~\mathrm{ms}\). Each packet contains a
constant identifier, a sequence number, a sample index, and one value from a
small test data set. This provides a repeatable 2.4~GHz radio transmission
whose packets can be used to explore reception and signal observation.

The Pluto SDR scripts also use a continuous-wave (CW) signal as a simple
frequency reference. In the laboratory setup, the transmitter and receiver
local oscillators are set to \(433~\mathrm{MHz}\), while the transmitter's
digital tone is offset by \(100~\mathrm{kHz}\). The resulting signal is
therefore expected near \(433.1~\mathrm{MHz}\). With a sampling rate and
receiver bandwidth of \(2~\mathrm{MHz}\), the captured I/Q buffer can be
transformed using a fast Fourier transform (FFT). The frequency of the largest
spectral bin gives an estimate of the received tone's offset from the receiver
local oscillator.

The second laboratory script extends this measurement with a live spectrum
display. Before calculating the FFT, it applies a Hann window to the samples
to reduce spectral leakage caused by observing a finite data buffer. The
frequency axis is shifted so that zero frequency is at the centre, then
translated by the receiver's local-oscillator frequency to display absolute
frequency. A marker indicates the largest spectral peak. These exercises
introduce the receiver settings and analysis techniques used in the project;
the \(433~\mathrm{MHz}\) CW test is a preliminary measurement and is distinct
from the project's 2.4~GHz transmission.

\subsection{2.4~GHz transmission and I/Q acquisition}

The ESP-NOW transmitter is paired with a Pluto SDR receiver implemented as a
GNU Radio flowgraph. It is tuned to \(2412~\mathrm{MHz}\) and samples at
\(20~\mathrm{MSample/s}\). The flowgraph passes the received samples to an
IEEE 802.11 physical-layer decoder, then examines the decoded frame for the
test packet identifier and expected data values. Sequence numbers also make
it possible to report gaps in the received packet stream. This experiment
demonstrates packet reception over the air, in addition to the raw spectrum
measurements of the CW exercises.

The project also includes a Pluto SDR CW transmitter configured with a
\(2.4~\mathrm{GHz}\) local oscillator and a \(100~\mathrm{kHz}\) digital tone,
placing the emitted tone at approximately \(2.4001~\mathrm{GHz}\). A separate
basic receiver script tunes to \(2.4~\mathrm{GHz}\), samples at
\(2~\mathrm{MSample/s}\), and prints a captured I/Q buffer. This is a raw
acquisition check; it does not attempt to decode a message.

\paragraph{BPSK waveform and recovery.}

For the project signal, the BPSK script configures the Pluto SDR transmitter
and receiver at a carrier frequency of \(2.4~\mathrm{GHz}\), with a sampling
rate of \(1~\mathrm{MSample/s}\). It generates a reproducible frame of
\(1024\) pseudo-random bits. Each bit is mapped to one of two antipodal
symbols, \(-1\) or \(+1\), and repeated for eight samples. The corresponding
symbol rate is
\[
    R_s = \frac{f_s}{N_s}
        = \frac{1~\mathrm{MSample/s}}{8}
        = 125~\mathrm{ksymbol/s},
\]
where \(f_s\) is the sampling rate and \(N_s\) is the number of samples per
symbol. The resulting complex-valued waveform is transmitted cyclically,
allowing the receiver to capture repeated copies of the same known frame.

The receiver acquires complex samples \(x[n]=I[n]+jQ[n]\). One initial
acquisition is discarded after configuration to avoid using a transition
buffer. The received samples are then processed to compensate for a constant
phase rotation, integrate over each symbol with an eight-sample rectangular
filter, and select the best of the eight possible sample timings. Because the
transmitted frame is known, correlation is used to find its alignment in the
received buffer and resolve the possible BPSK sign ambiguity. A hard decision
on the real part of each recovered symbol produces the estimated bit sequence.
The script compares it with the original bits and reports the bit-error rate
(BER); it also displays the recovered constellation and received spectrum.

\subsection{Phase 1 outcome and scope}

Together, the laboratory work and project scripts demonstrate how to generate
a known RF waveform, configure an SDR receiver, inspect the received spectrum,
and recover a digital sequence from I/Q samples. They also provide a
repeatable signal-processing chain that can be used to check transmission and
acquisition settings. The scripts report measurements such as the detected
frequency and BER when run, but no fixed numerical results are included here
because they depend on the hardware setup and the individual acquisition.

These experiments use a single receive channel for the BPSK recovery and
spectrum measurements. They therefore do not, on their own, estimate the
direction of arrival. Spatial phase comparison across multiple receive
channels, together with calibration, is required for that task and is
addressed in the following phase.